\documentclass[../DoAn.tex]{subfiles}
\begin{document}

\begin{center}
    \Large{\textbf{TÓM TẮT NỘI DUNG ĐỒ ÁN}}\\
\end{center}
\vspace{1cm}

Đồ án nghiên cứu và xây dựng hệ thống trợ lý học vụ dựa trên kiến trúc
\textit{Retrieval-Augmented Generation} (RAG) cho sinh viên Đại học Bách khoa
Hà Nội. Hệ thống hướng tới việc hỗ trợ tra cứu bằng ngôn ngữ tự nhiên trên các
nguồn tri thức học vụ nội bộ, đồng thời giảm rủi ro trả lời thiếu căn cứ bằng
cách gắn quá trình sinh câu trả lời với tài liệu đã truy hồi. Phạm vi tri thức
chính được tổ chức thành các nhóm dữ liệu về chương trình đào tạo
(\texttt{ctdt}), quy định học vụ (\texttt{quydinh}), kế hoạch và thông báo
(\texttt{kehoach}), tài liệu hỗ trợ sinh viên (\texttt{stsv}); ngoài ra còn có
vùng dữ liệu phục vụ kiểm thử và tải lên tài liệu mới.

Phiên bản hệ thống được hiện thực trong mã nguồn theo kiến trúc nhiều lớp. Lớp
giao tiếp sử dụng FastAPI để cung cấp các API chat thường, chat có metadata theo
dõi luồng xử lý, chat dạng streaming, tìm kiếm kiểm tra truy hồi, quản lý phiên
hội thoại, xác thực người dùng, quản trị tài liệu, bookmark, feedback, lookup và
notification. Lớp điều phối trung tâm là \texttt{RAGPipeline}: câu hỏi đầu vào
được phân loại độ phức tạp để tách giữa hội thoại xã giao, câu hỏi RAG đơn giản
và câu hỏi phức tạp. Với câu hỏi đơn giản, pipeline thực hiện luồng RAG cổ điển;
với câu hỏi nhiều nguồn hoặc so sánh, hệ thống có thể phân rã thành các truy vấn
con theo miền tri thức; với câu hỏi phức tạp khác, hệ thống dùng agent dựa trên
LangGraph và vẫn có cơ chế quay về RAG khi agent không sẵn sàng hoặc xử lý lỗi.

Trước khi truy hồi, lớp xử lý truy vấn kết hợp nhiều thành phần: bộ định tuyến
độ phức tạp, bộ phân loại ý định và miền dữ liệu, bộ chọn collection, bộ phản
chiếu truy vấn và bộ phân rã truy vấn. Câu hỏi có thể được viết lại thành truy
vấn độc lập dựa trên lịch sử hội thoại và ngữ cảnh người dùng khi thật sự cần
thiết; đồng thời hệ thống trích xuất các thực thể học vụ như mã ngành, khóa sinh
viên, học kỳ hoặc mã học phần để tạo bộ lọc phù hợp. Mã nguồn cũng đặt các ràng
buộc nhằm hạn chế việc ngữ cảnh hồ sơ hoặc lịch sử hội thoại bị chèn sai vào
những câu hỏi chung, đặc biệt là các câu hỏi yêu cầu thông tin mới nhất.

Khối truy hồi được thiết kế theo hướng lai để kết hợp tìm kiếm ngữ nghĩa và tìm
kiếm từ khóa. Dịch vụ truy hồi dùng hai biểu diễn embedding BGE-M3 và E5
multilingual cho vector search trên Qdrant, đồng thời dùng Elasticsearch cho
keyword search và hỗ trợ lọc metadata. Kết quả từ nhiều collection được chuẩn
hóa điểm, hợp nhất theo trọng số, loại trùng và xếp hạng lại bằng BGE reranker.
Sau truy hồi, hệ thống còn áp dụng bộ lọc hiệu lực tài liệu để tránh dùng văn
bản đã bị thay thế khi có thông tin lineage phù hợp, và bộ phân giải tham chiếu
để bổ sung các đoạn liên quan đến các dẫn chiếu pháp lý như điều, khoản trong
cùng văn bản. Đối với câu hỏi về kế hoạch, hạn đăng ký hoặc thông báo có tính
thời điểm, pipeline có thể ưu tiên miền \texttt{kehoach}; khi dữ liệu nội bộ
không đủ hoặc có nhu cầu thông tin mới, cơ chế Tavily web fallback được cấu hình
để tìm kiếm có kiểm soát trên các nguồn chính thống.

Đồ án không chỉ xử lý luồng hỏi đáp mà còn xây dựng vòng đời dữ liệu phục vụ RAG.
Pipeline quản trị tài liệu nhận PDF từ admin, lưu tệp gốc, chuyển đổi sang
Markdown, làm sạch nội dung, chia chunk, tạo embedding và lập chỉ mục đồng thời
vào Qdrant và Elasticsearch. Trạng thái tài liệu được theo dõi theo từng bước
từ tải lên đến indexed trong MongoDB; các chunk trung gian được lưu để phục vụ
kiểm tra, phê duyệt và rollback. Bên cạnh luồng tải lên qua API, các script crawl
và index ngoại tuyến hỗ trợ cập nhật kho dữ liệu miền học vụ, bảo trì metadata và
quản lý dòng đời văn bản.

Ở mức hệ thống, MongoDB lưu người dùng, phiên hội thoại, lượt hỏi đáp, trace của
agent, tài liệu, chunk, bookmark, feedback và notification. Redis là thành phần
tùy chọn cho session, history cache, response cache và rate limiting theo hướng
fail-soft để backend vẫn hoạt động khi cache không khả dụng. Cơ chế xác thực hỗ
trợ JWT, đăng nhập Microsoft OAuth và tài khoản thủ công; các thao tác quản trị
tài liệu được bảo vệ bằng kiểm tra quyền. Phần giao diện gồm web app React/Vite,
ứng dụng mobile Expo/React Native và package TypeScript dùng chung để đồng bộ
API contract giữa các client.

Để kiểm soát chất lượng, mã nguồn cung cấp lớp đánh giá ngoại tuyến và kiểm thử
hồi quy. Bộ \texttt{evaluation} tách kiểm tra truy hồi theo chính sách hiện hành
khỏi đánh giá hội thoại trên dữ liệu lịch sử, đồng thời hỗ trợ lưu artifact,
dashboard metrics và benchmark chiến lược tìm kiếm. Các kiểm thử bao phủ routing,
reflection, retrieval fusion, reference resolution, API contract, upload
pipeline, agent và các luồng cache liên quan. Nhờ đó, nội dung đồ án không chỉ
trình bày một chatbot sinh văn bản, mà mô tả đầy đủ một hệ thống RAG học vụ có
khả năng nhập dữ liệu, truy hồi lai, suy luận theo độ phức tạp truy vấn, theo dõi
nguồn trả lời, vận hành đa client và đánh giá chất lượng theo thời gian.

\begin{flushright}
Sinh viên thực hiện\\
\begin{tabular}{@{}c@{}}
\textit{(Ký và ghi rõ họ tên)}
\end{tabular}
\end{flushright}

\end{document}
